\section{Introducción}

\textbf{Nombre provisorio:} Take a Break.

Take a Break es una app para Android que convierte el tiempo que una persona pasa lejos del celular en puntos, y esos puntos se canjean por premios en comercios del barrio. No hay que activar nada. La app se entera sola de que el teléfono estuvo quieto y, si la pausa cumple las reglas, suma.

\begin{nota}
\centering
\textbf{Problema} $\rightarrow$ \textbf{Usuario} $\rightarrow$ \textbf{Solución} $\rightarrow$ \textbf{Valor}\\[4pt]
Las herramientas para usar menos el celular miden o bloquean, pero no dan nada a cambio
$\rightarrow$ adultos jóvenes que ya lo intentaron y abandonaron
$\rightarrow$ detección pasiva de pausas y premios reales cerca
$\rightarrow$ un hábito más sano para el usuario y clientes nuevos para el comercio.
\end{nota}

\subsection{Cómo está organizado este documento}

Trabajamos con design thinking, así que el documento sigue sus cinco etapas. La tabla muestra dónde está cada punto que pide la consigna.

\begin{table}[H]
\centering\small
\begin{tabularx}{\textwidth}{@{}l L l@{}}
\toprule
\textbf{Etapa} & \textbf{Puntos del entregable} & \textbf{Sección} \\
\midrule
Empatizar  & 2. Descripción del problema · 3. Usuarios · 4. Contexto de uso & \ref{sec:empatizar} \\
Definir    & Punto de vista, preguntas guía y principios de diseño & \ref{sec:definir} \\
Idear      & 1. Nombre · 5. Propuesta · 6. Por qué móvil · 7. Caso de negocio · 8. Alcance & \ref{sec:idear} \\
Requisitos & 9. Requisitos funcionales · 10. Requisitos no funcionales & \ref{sec:requisitos} \\
Prototipar & 11. Diseño en Figma · 12. Flujo de pantallas & \ref{sec:prototipar} \\
Testear    & Plan de validación con usuarios & \ref{sec:testear} \\
Construir  & 13. Arquitectura · 14. Offline first · 15. Tecnologías & \ref{sec:arquitectura} a \ref{sec:tecnologias} \\
Equipo     & 16. Repositorio · 17. Integrantes y roles & \ref{sec:equipo} \\
\bottomrule
\end{tabularx}
\end{table}

Los archivos de diseño están en el repositorio \texttt{docs}, carpeta \texttt{figma/}. Todas las pantallas aparecen en el Anexo~\ref{anx:pantallas}.
